% Companion to the checked explicit wild cover.  Independent audit requested.
\begin{proposition}[Proved: the source resolution of the explicit fierce double cover]
\label{D:wild-genus2-source}
Let $T\to S$ be the cover in Example~\ref{D:wild-genus2-single-curve}.
The surface $T$ has a projective resolution $V\to T$ whose
exceptional fibre is one smooth genus-two curve $C'$, with
$(C')^2=-2$.  There is a finite morphism $V\to Y$ to the
resolution of $S$ in that example.  Its restriction $C'\to C$
is the relative Frobenius, after the displayed identifications
of both curves with $w^2+w=z^5$.  The discrepancy of $C'$
over the Gorenstein surface $T$ is $-2$.
\end{proposition}
\begin{proof}
Use weights $(1,3,1)$ on $(c,d,t)$ and the weighted Rees
construction of Example~\ref{D:wild-genus2-single-curve}.
The algebra is generated by $cs,ds,ds^2,ds^3,ts$ over
$k[c,d,t]$, and its relative Proj is proper and is an
isomorphism away from the origin.  It is projective: a
sufficiently divisible Veronese is generated in degree one,
so the same relative Proj is a closed subscheme of a
projective space over the affine base; see
\cite[Lemma 10.56.2, tag 0EGH]{StacksBWildSource}.
Let $V$ be the strict
transform of $T$.  Its reduced exceptional curve is
\[
                 d^2+t^3d+tc^5=0
                       \quad\hbox{in }\mathbb P(1,3,1).
\]
On the $t$ chart set $c=tz$ and $d=t^3w$.  Dividing the
equation of $T$ by $t^6$ gives
\begin{equation}\label{D:source-t-chart}
 w^2+(1+tz^4+t^3z^2)w+
            z^5+t^3z^7+t^4z^6+t^5z^5=0.
\end{equation}
Its exceptional curve is $w^2+w=z^5$, and the $w$
derivative is one along this curve.  Thus this chart is
smooth on a neighborhood of the exceptional curve.

The only point of the exceptional curve outside this chart
has $t=d=0$ and $c\ne0$ in the weighted projective coordinates.
Since $c$ has weight one, its chart is an ordinary affine
chart.  Put
\[
                   c=r,\qquad t=r\tau,\qquad d=r^3\delta.
\]
The strict-transform equation is
\begin{equation}\label{D:source-c-chart}
 \delta^2+(\tau^3+r+r^3\tau^4)\delta
       +\tau+r^3\tau^2+r^4\tau^4+r^5\tau^6=0.
\end{equation}
At the exceptional point $(r,\tau,\delta)=(0,0,0)$ the
$\tau$ derivative is one.  The two charts therefore prove
that $V$ is smooth near its exceptional fibre.  Off that
fibre the morphism is an isomorphism over the smooth
punctured surface $T$.  They also show that its only
exceptional curve $C'$ is the smooth proper completion of
$w^2+w=z^5$, of genus two by the double-cover computation
in the preceding example.

The original function $c$ equals $tz$ in the first chart.
Its strict transform meets $C'$ at precisely the two points
\[
                    (z,w)=(0,0),\qquad (z,w)=(0,1).
\]
Each intersection is transversal: the $w$ derivative in
\eqref{D:source-t-chart} is a unit there, so $t,z$ are
regular parameters.  In the second chart $c=r$ is only
the exceptional equation, so there is no additional
intersection at infinity.  Intersecting the principal
divisor of $c$ with $C'$ yields
\[
                            0=(C')^2+2.
\]

We give the map to $Y$ on the same charts.  On the $t$
chart of the target it is
\[
                    x=z^2+tz,\qquad y=w+z^5,\qquad t=t.
\]
The equation of $Y$ pulls back to
\eqref{D:source-t-chart}.  On a neighborhood of the source
point at infinity put
\[
             A=1+r\tau^2,\qquad B=1+\tau^2\delta.
\]
Both are units there.  The target's infinity coordinates
are
\[
             \rho=\frac{rB}{A^2},\qquad
             \xi=\frac{A^5}{B^2},\qquad
             \tau_Y=\frac{\tau A^2}{B}.
\]
These formulas follow from $a=r^2A$, $b=r^5B$ and $t=r\tau$,
and agree on overlaps.  Together with the original map off
the exceptional fibre they define a morphism $V\to Y$.
It is proper because $V$ is proper over $S$ and $Y$ is
separated over $S$.  Its fibres are finite away from $C$.
On $C'$ the map is
\[
                      (z,w)\longmapsto(z^2,w^2),
\]
since $w+z^5=w^2$ there.  It is consequently finite on
the exceptional fibre as well.  Properness and
quasi-finiteness make $V\to Y$ finite.  Thus $V$ is also
the normalization of $Y$ in the quadratic extension.

Finally $T$ is Gorenstein, and adjunction gives
$K_V\cdot C'=2g(C')-2-(C')^2=4$.  If $a$ is the
discrepancy of its sole exceptional curve, then
$a(C')^2=4$, whence $a=-2$.
\end{proof}

\begin{corollary}[Proved: the next quadratic step and the local pro-two quotient]
\label{D:wild-genus2-pro2}
Every finite normal cover of $T$ of degree two that is
\'etale off its singular point is tame at $C'$.
Let $U_T$ and $U_S$ denote the punctured spectra of the
strict henselizations at the singular points of $T$ and
$S$, respectively.  Then
\[
                 \pi_1(U_T)^{(2)}=1,\qquad
                 \pi_1(U_S)^{(2)}=\mathbb Z/2\mathbb Z,
\]
where the superscript denotes the maximal pro-two quotient.
This does not assert that the full wild inertia of either
punctured fundamental group is finite.
\end{corollary}
\begin{proof}
Apply Proposition~\ref{D:single-curve-residue-tameness} with
$g=2$, $d=2$ and $n=2$.  If any extension of the exceptional
valuation has inseparable residue, its residue degree is
at least two.  The degree formula then forces a unique
extension, with $e=1$, $f=2$ and $f^{\mathrm{sep}}=1$.
The inequality in that proposition gives
\[
                              2\delta\leq1.
\]
Its equality criterion for tame ramification instead gives
$\delta\geq1$ for an inseparable residue extension, a
contradiction.  Thus all residue extensions are separable,
and the same proposition gives tameness.  This argument
makes no assumption on the number of components in a
resolution of the covering surface.

We explain the strict-local assertion to keep its scope
separate from completion or algebraization claims.  A
finite cover of the punctured strict henselization descends,
by finite presentation, to an \'etale neighborhood
of the singular point.  Normalize that neighborhood and
shrink it if necessary.  The preceding argument applies
there, since the pulled-back resolution has the same
exceptional curve and normal degree.  A degree-two cover
of $U_T$ is therefore tame at $C'$.  In characteristic two
its ramification indices are prime to two and divide two,
so they are all one.  The normalization of the regular
resolution in this cover is unramified at every
codimension-one point, and hence is finite \'etale
by purity of the branch locus
\cite[Lemma 58.21.4, tag 0BMB]{StacksBWildSource}.
Proper henselian
comparison identifies its finite \'etale category
with that of $C'$; see
\cite[Lemma 58.9.1, tag 0A48]{StacksBCuspLocal}.
Here we pass from the scheme-theoretic closed fibre to
its reduction $C'$ using invariance of finite \'etale
covers under nilpotent thickenings.

There is no nontrivial \(\mathbb Z/2\mathbb Z\)-torsor on
$C'$.  Indeed, in the two-chart \v Cech basis
$h_1=w/z$, $h_2=w/z^2$ of $H^1(C',\mathcal O_{C'})$,
the equation $w^2+w=z^5$ gives
\[
                              F(h_1)=h_2,\qquad F(h_2)=0.
\]
The extra terms $z^3,z$ are regular on the finite chart,
and $w/z^4$ is regular on the infinity chart.  Thus $F$
is nilpotent, and the Artin--Schreier exact sequence gives
$H^1_{\text{ét}}(C',\mathbb F_2)=0$.  We conclude
that $\pi_1(U_T)$ has no quotient of order two.  Every
nontrivial finite two-group has such a quotient, so its
maximal pro-two quotient is trivial.

The given double cover remains connected on strict
henselization: its finite normalization is normal, its
closed fibre has only one point, and a finite algebra over
a henselian local ring is a product of its local factors.
This local normalization is strictly henselian and is
ind-\'etale over $\mathcal O_{T,o}$; it is therefore
the strict henselization defining $U_T$.
Thus $\pi_1(U_T)$ is the normal subgroup of index two in
$\pi_1(U_S)$ defined by that cover.  In any finite
two-group quotient of $\pi_1(U_S)$ the image of this
subgroup must be trivial, since $\pi_1(U_T)$ has no
nontrivial two-group quotient.  All such quotients
therefore factor through the given group of order two,
which proves the second assertion.
\end{proof}

\begin{remark}[Exact finite verification]
The companion program
\texttt{check\_wild\_genus2\_source.py} verifies 22 exact
identities over $\mathbb F_2$, including both resolution
charts and their overlap, the finite map on both charts,
the two intersections used for $(C')^2$, and the
Frobenius action on the two cohomology classes.  Its
recorded runtime is about $0.14$ seconds.  The arguments
for properness, connectedness, purity and the profinite
conclusions are the mathematical arguments above.
\end{remark}

\begin{remark}[Remaining scope of the example]
The source resolution and the pro-two calculation rule out
repeating this particular wild quadratic step.  They do
not rule out wild inertia inside a finite group that is
not a two-group, or after passing first to other finite
punctured covers.  A bound for those images, and the
general D.1/D.2(ii) finite-generation statements, do not
follow from this calculation.
\end{remark}
